\subsection{PRIMITIVES OF RATIONAL FUNCTIONS}

Formula (31) allows us to evaluate the primitive of an arbitrary rational function \(r(x)\) of a real variable x, with real or complex coefficients. We know (A.VII.7) that such a function can be written (in unique manner) as a finite sum of terms, which are:

\begin{enumerate}
\def\labelenumi{\alph{enumi})}
\item
  either of the form \(ax^{p}\) (p an integer \(\geqslant 0\) , a a complex number);
\item
  or of the form \(a/(x - b)^{m}\) (m an integer \(\geqslant 0\) , a and b complex numbers).
\end{enumerate}

Now it is easy to obtain a primitive of each of these terms:

\begin{enumerate}
\def\labelenumi{\alph{enumi})}
\tightlist
\item
  a primitive of \(ax^p\) is \(a\frac{x^{p + 1}}{p + 1}\) ;
\end{enumerate}

\(b)\) if \(m > 1\) a primitive of \(a / (x - b)^{m}\) is \(\frac{a}{(1 - m)(x - b)^{m - 1}}\) ;

\begin{enumerate}
\def\labelenumi{\alph{enumi})}
\setcounter{enumi}{2}
\tightlist
\item
  finally, from formulae (10) (III, p.~93) and (31) (III, p.~101), a primitive of \(\frac{a}{x - b}\) is \(a\log |x - b|\) if \(b\) is real, \(a\log (x - b)\) if \(b\) is complex. In the last case, if \(b = p + iq\) , one has furthermore (III, p.~100, formulae (30))
\end{enumerate}

\[
\log (x - b) = \log \sqrt {(x - p) ^ {2} + q ^ {q}} + i \operatorname{Arc} \tan \frac {x - p}{q} \pm i \frac {\pi}{2}.
\]

We postpone the examination of more practical methods for explicitly determining a primitive of a rational function given explicitly to the part of this work dedicated to Numerical Calculus.

One can reduce to the evaluation of a primitive of a rational function:

\(1^{\circ}\) the evaluation of a primitive of a function of the form \(r(e^{ax})\) , where \(r\) is a rational function and \(a\) a real number; indeed by the change of variable \(u = e^{ax}\) one reduces to finding a primitive for \(r(u) / u\) ;

\(2^{\circ}\) the evaluation of a primitive of a function of the form \(f(\sin ax, \cos ax)\) , where \(f\) is a rational function of two variables and \(a\) is a real number; the change of variables \(u = \tan(ax/2)\) reduces this to finding a primitive for

\[
{\frac {2}{1 + u ^ {2}}} f \left({\frac {2 u}{1 + u ^ {2}}}, {\frac {1 - u ^ {2}}{1 + u ^ {2}}}\right).
\]

